\begin{statementbox}
\begin{center}\begin{tikzpicture}[font=\small,>=stealth,line join=round]
\draw[gray!40,rounded corners=3pt] (0,0) rectangle (14.2,8.3);\node at (7.1,7.700000000000001) {左右等距、函数值同高：对称轴就在中间};
\begin{scope}[shift={(1.7923076923076926,2.038732394366197)},x=1.153846153846154cm,y=0.647887323943662cm]\draw[->,gray] (-0.6,0)--(4.6,0) node[right]{$x$};\draw[->,gray] (0,-0.6)--(0,6.5) node[above]{$y$};\draw[blue,thick,solid,domain=-0.1:4.1,samples=180] plot(\x,{(\x-2)^2+1});\fill[orange!85!black] (0,5) circle(2.2pt);\fill[orange!85!black] (1,2) circle(2.2pt);\fill[orange!85!black] (2,1) circle(2.2pt);\fill[orange!85!black] (3,2) circle(2.2pt);\fill[orange!85!black] (4,5) circle(2.2pt);\draw[violet,dashed] (2,-.5)--(2,6);\node[above,violet] at (2,6){$x=2$};\draw[orange!85!black,dashed] (0,5)--(4,5);\draw[orange!85!black,dashed] (1,2)--(3,2);\foreach \v in {0,1,2,3,4}{\draw[gray] (\v,0)--(\v,-.1);\node[below=5pt] at (\v,-.1){$\v$};}\end{scope}\node[anchor=west,align=left,text width=5.0cm] at (8.4,4.0) {$f=(x-2)^2+1$\\[9pt]$f(0)=f(4)=5$\\[9pt]$f(1)=f(3)=2$\\[9pt]每对横坐标之和为 $4$\\[9pt]轴的位置：$4\div2=2$};
\node[text width=12.6cm,align=center] at (7.1,.65) {图示具体配对；完整证明要求定义域左右成对，且每一对都同高。};
\end{tikzpicture}\end{center}


  \textbf{一句话记住：轴对称，横坐标取平均，纵坐标相等。}
  这里说的是关于\textbf{竖直直线}的对称：两点横坐标的平均数给出轴的位置；两点高度相同。

  \textbf{\textcolor{CStatement}{先想一想：这条结论解决什么问题？}}
  有时函数的图象不好画，却要判断它是否能沿某条竖直线对折重合；
  有时只知道图象右半边，却要算左半边的函数值。
  这时不必先把整张图画出来，只要检查对称轴两侧的点是不是“一样高”。

  \medskip
  \textbf{\textcolor{CStatement}{条件（左右两边都必须有点）}}
  设函数为 $y=f(x)$，准备检查的对称轴是 $x=a$。
  只要 $a+t$ 可以代入函数，$a-t$ 也必须可以代入；反过来一样。
  也就是说，定义域本身要以 $a$ 为中间位置，左右成对。

  \medskip
  \end{statementbox}\clearpage\begin{statementbox}\textbf{\textcolor{CStatement}{核心结论（左右一样远，函数值一样）}}
  在上述条件下，图象关于直线 $x=a$ 对称，\textbf{当且仅当}
  对每个允许的 $t$，都有
  \[
    \boxed{f(a+t)=f(a-t)}.
  \]
  \textbf{读成一句话：}从 $a$ 往右走 $t$，和从 $a$ 往左走 $t$，到达的高度相等。
  这里 $t$ 是带方向的水平偏移，实际距离为 $|t|$；$t$ 可正可负。

  \medskip
  \textbf{\textcolor{CStatement}{最适合做题的写法}}
  把左边位置记为 $x$，右边与它配对的位置就是 $2a-x$，所以
  \[
    \boxed{f(x)=f(2a-x)\quad(x\text{ 在定义域内})}.
  \]
  两个位置的\textbf{平均数是 $a$，和是 $2a$}：
  \[
    \underbrace{\frac{x+(2a-x)}{2}}_{\text{中间位置}}=a.
  \]
  使用时先保证这两个位置都在定义域内。

  \medskip
  \textbf{\textcolor{CStatement}{举个一眼能看懂的例子}}
  对称轴 $x=2$：$1$ 与 $3$ 配对，$0$ 与 $4$ 配对，$-1$ 与 $5$ 配对。
  因此若图象关于 $x=2$ 对称，就一定有
  \[
    f(1)=f(3),\qquad f(0)=f(4),\qquad f(-1)=f(5).
  \]
  反过来，\textbf{仅验证这三对还不够}；必须能保证定义域内\textbf{所有}配对都成立。
\end{statementbox}
\clearpage
\begin{statementbox}
\textbf{同一个判定的常见等价写法：先看实际输入，再找中间位置}
\begin{center}\begin{tikzpicture}[font=\small,>=stealth,line join=round]
\draw[gray!40,rounded corners=3pt] (0,0) rectangle (14.2,8.5);\node at (7.1,7.9) {例6：真实代入位置相加恒为6，中点才是轴};
\draw[->,gray] (1,2)--(13,2) node[right]{$x$};\draw[violet,dashed] (7,1.5)--(7,6.7);\draw[blue,thick] (3,4)--(11,4);\fill[orange!85!black] (3,4) circle(2.4pt);\draw[dashed,gray] (3,2)--(3,4);\node[below=7pt] at (3,2){$t+1$};\fill[orange!85!black] (11,4) circle(2.4pt);\draw[dashed,gray] (11,2)--(11,4);\node[below=7pt] at (11,2){$5-t$};\node[above,violet] at (7,6.7){$x=3$};\node[above=8pt] at (10,4){同高};
\node[text width=12.6cm,align=center] at (7.1,.65) {横向以 $t=0$ 的配对 $1,5$ 示意；一般位置之和恒为6，不指定函数值。};
\end{tikzpicture}\end{center}


以下先取定义域为 $\mathbb R$，$a,p,q$ 为固定实数，每行等式均对所有实数自变量成立：
\begin{center}
\renewcommand{\arraystretch}{1.4}
\begin{tabular}{p{.56\linewidth}|p{.28\linewidth}}
等式的写法 & 能读出的一条对称轴\\\hline
$f(a+t)=f(a-t)$ & $x=a$\\
$f(x)=f(2a-x)$ & $x=a$\\
$f(p+t)=f(q-t)$ & $x=\dfrac{p+q}{2}$\\
$f(p-t)=f(q+t)$ & $x=\dfrac{p+q}{2}$
\end{tabular}
\end{center}
\textbf{只需一个读式规则：}例如 $(p+t)+(q-t)=p+q$，两个输入的平均数恒为 $(p+q)/2$；对应函数值又相等，因此可读出这条竖直对称轴。令 $u=p+t$，等式就是 $f(u)=f(p+q-u)$，仍是原来的判定。最后一行只需将 $t$ 换为 $-t$。

图中 $p=1,q=5$，两输入是 $1+t,5-t$，所以中间位置为 $3$；图仅展示 $t=0$ 时的一对等高点。
\textbf{其他定义域：}须保证每个 $u\in D$ 都有 $p+q-u\in D$，且等式覆盖全部这些配对。只检查少数 $t$，不能判定整幅图象对称。
\end{statementbox}
